\chapter*{Tóm tắt}

Khóa luận này nghiên cứu bài toán phân đoạn tổn thương xương trên ảnh X-quang theo hướng có tương tác với người dùng. Nói đơn giản hơn, thay vì để mô hình tự tìm hết mọi tổn thương trên ảnh, người dùng sẽ khoanh trước một hộp giới hạn sơ bộ rồi mô hình mới phân đoạn chi tiết hơn trong vùng đó. Cách làm này hợp với ảnh X-quang xương hơn vì tổn thương thường nhỏ, biên không rõ hoặc bị chồng với các cấu trúc giải phẫu khác.

Từ hướng làm đó, khóa luận đề xuất PGA-UNet (\textit{Prompt-Guided Attention U-Net}), là một mô hình CNN gọn nhẹ với khoảng 3 triệu tham số. Hộp giới hạn được đổi thành một bản đồ nhiệt câu nhắc dạng cao nguyên, trong đó phần biên được làm mềm bằng bộ lọc Gaussian. Tín hiệu câu nhắc này được đưa vào bộ mã hóa qua Cổng không gian câu nhắc (PSG) và được dùng tiếp ở nhánh giải mã qua Cơ chế chú ý có điều kiện (CAD).

Mô hình được thực nghiệm trên hai bộ dữ liệu là BTXRD và FracAtlas. Phần đánh giá gồm so sánh với các mô hình cơ sở, nghiên cứu loại bỏ thành phần, đánh giá chéo và thử nghiệm với các câu nhắc bị sai lệch. Cũng cần nói trước là các kết quả trong khóa luận đều nằm trong thiết lập phân đoạn có hướng dẫn bằng hộp giới hạn, trong đó câu nhắc được tạo từ nhãn gốc để mô phỏng thao tác người dùng. Nói cách khác, khi so với U-Net và Attention U-Net thì kết quả chủ yếu phản ánh khác biệt giữa mô hình có câu nhắc và mô hình không có câu nhắc. Còn phép so sánh cùng điều kiện đầu vào thì chủ yếu được làm với SAM-Med2D ở độ phân giải $256\times256$.

Trong thiết lập này, trên BTXRD ở độ phân giải $512\times512$, PGA-UNet đạt Dice $0{,}8607$, còn U-Net chỉ đạt $0{,}4740$. Khi câu nhắc bị lệch tâm, mô hình vẫn giữ được Dice $0{,}8423$. Với phép so sánh cùng điều kiện câu nhắc với SAM-Med2D tại độ phân giải $256\times256$, PGA-UNet đạt Dice $0{,}8433$, cao hơn mức $0{,}7541$ của SAM-Med2D đã tinh chỉnh. Trên FracAtlas, mô hình đạt Dice $0{,}8169$ và xu hướng chung cũng gần như vậy. Phần ablation cũng cho thấy PSG, CAD và bản đồ nhiệt Gaussian khi đi cùng nhau giúp mô hình giữ kết quả ổn hơn khi câu nhắc bị sai lệch.

Ngoài ra, khóa luận còn đo hiệu quả tính toán trong cùng một môi trường thực nghiệm. Ở độ phân giải $256\times256$, PGA-UNet nhỏ hơn SAM-Med2D khoảng 92 lần về số tham số, 11,9 lần về FLOPs và có độ trễ suy luận thấp hơn trên cả GPU lẫn CPU. Khóa luận cũng thử thêm một luồng xử lý hai giai đoạn bằng cách đặt mô-đun sàng lọc EfficientNet-B3 trước PGA-UNet. Phần này cho thấy nếu bước sàng lọc bỏ sót ảnh bệnh lý thì hiệu năng toàn hệ thống giảm đi rõ rệt.

\noindent\textbf{Từ khóa:} Phân đoạn ảnh y khoa, phân đoạn tương tác, câu nhắc trực quan, ảnh X-quang xương, PGA-UNet.
